% Ad Atticum 14.10 --- headnote
% ad-atticum-14-10, 19 April 44 BC, in Cumano

\textit{Cicero to Atticus, written at the Cumean villa on 19
April 44 BC --- Perseus dateline \textit{Scr. in Cumano xiii K.
Mai. a. 710 (44)}. The letter opens in a high pitch of
retrospective anger: Brutus has retired to Lanuvium, Trebonius
is slipping off to Asia by side-roads, and all Caesar's
unpublished \textit{acta} are being honoured as binding. Cicero
revisits his own cry from the Capitol on the Ides --- the
praetors should have summoned the Senate then and there --- and
his warning that the cause was lost the moment a public funeral
was granted. The catalogue of beneficiaries (Tebassi, Scaevae,
Fangones, Curtilius, Censorinus, Messalla, Plancus, Postumus)
is a roll-call of Caesarian planters whose tenure depends on
the regime's survival.}

\textit{Two Greek tags frame the mood. ``\textit{g\=en pro
g\=es}'' --- ``one land after another'' --- is the proverb of
the exile (Aeschylus, \textit{Prometheus} 682); Cicero is
contemplating flight. Atticus's flight is more ethereal:
``\textit{hyp\=enemios}'', ``wind-borne'', empty. The third
Greek word, daggered in the manuscripts (\textit{rhixothemin}),
is a corruption beyond clean restoration; the English keeps
the obelus visible. The closing items are quieter: the young
Octavian has arrived at Naples and is going to accept his
inheritance; the Cluvian estate (the recent bequest) seems to
be approaching its target yield; and Quintus the elder is once
again writing furiously about his son.}
